\documentclass[11pt,a4paper]{article}
\usepackage[T1]{fontenc}
\usepackage{lmodern,microtype,amsmath,amssymb,amsthm,booktabs}
\usepackage[margin=23mm]{geometry}
\usepackage[colorlinks=true,linkcolor=blue,urlcolor=blue]{hyperref}
\newtheorem{theorem}{Theorem}
\newtheorem{lemma}[theorem]{Lemma}
\newtheorem{corollary}[theorem]{Corollary}
\newcommand{\E}{\mathbb E}
\newcommand{\Prb}{\mathbb P}
\newcommand{\Law}{\operatorname{Law}}
\newcommand{\TV}{\mathrm{TV}}
\newcommand{\ii}{\imath}
\title{A Moment Classification of Exact Dyadic Causal Overhead\\[3pt]
\large Bounded, logarithmic, and square-root regimes at equal entropy}
\author{A constructive note for the history-indexed C3-C model}
\date{26 September 2026}
\begin{document}
\maketitle
\begin{abstract}
We classify the order of the exact causal seed overhead for all pairs of
finite pure dyadic laws with equal entropy. Unequal information variances
give order $\sqrt n$; equal variances but unequal third centered information
moments give order $\log n$; matching both moments gives bounded overhead,
with a positive limit unless the laws are permutations of one another.
An explicit pair on $79$ and $80$ symbols realizes the logarithmic regime.
The proof combines an exact optimal CDF recursion, a Fourier lower bound
on every entropy increment, and a smoothed envelope contraction. The lower
bounds also apply to arbitrary finite action laws with equal entropy.
\end{abstract}

\section{The classification}
All information quantities use base-two logarithms. For a finite law $A$,
write $\ii_A(x)=-\log_2 A(x)$ on its positive support. Define
\[
h_A=\E\ii_A(X),\quad
v_A=\E(\ii_A(X)-h_A)^2,\quad
\kappa_A=\E(\ii_A(X)-h_A)^3,\qquad X\sim A.
\]
A \emph{pure dyadic law} has each positive atom equal to $2^{-k}$ for an
integer $k\ge0$. The restriction is to individual powers of two, not to
arbitrary rational numbers with a power-of-two denominator.

In the history-indexed model C3-C, one initial discrete seed determines
all responses. At every positive-probability history, and under every
adaptive action policy, the next output must have exactly the selected
law $P$ or $Q$. The policy observes past actions and outputs, not the seed.
Let $C_n^{\mathrm{causal}}(P,Q)$ be the infimum seed entropy at horizon $n$.
The seed is allowed to depend on $n$.

\begin{theorem}[Complete equal-entropy classification for pure dyadic laws]
\label{thm:classification}
Let $P,Q$ be finite pure dyadic laws with $h_P=h_Q=h$, and set
$\Delta_n=C_n^{\mathrm{causal}}(P,Q)-nh$.
The optimum is attained by an explicit finite seed. The following cases
are exhaustive:
\[
\boxed{
\begin{array}{ll}
v_P\ne v_Q &: \Delta_n=\Theta(\sqrt n),\\[2pt]
v_P=v_Q,\ \kappa_P\ne\kappa_Q &: \Delta_n=\Theta(\log(n+1)),\\[2pt]
v_P=v_Q,\ \kappa_P=\kappa_Q &: \Delta_n=O(1).
\end{array}}
\]
In the bounded case, $\Delta_n$ is nondecreasing and converges to a finite
limit. This limit is positive unless $P,Q$ are permutations of one another;
in that exceptional case $\Delta_n=0$ for all $n$.
All asymptotic constants may depend on the fixed pair.
\end{theorem}
Thus, within this class, the second and third centered information
moments completely determine the growth order. Higher moments may affect
the bounded limit but do not create further unbounded orders.

\begin{corollary}[Necessary conditions without dyadic assumptions]
\label{cor:general}
For arbitrary finite action laws of equal entropy, unequal information
variances imply $\Delta_n=\Omega(\sqrt n)$. Equal variances and unequal
third centered information moments imply $\Delta_n=\Omega(\log(n+1))$.
In particular, bounded overhead requires matching both moments.
\end{corollary}

\section{Exact profile recursion and dyadic realization}
For any discrete law $A$, let $F_A(t)=\Prb\{\ii_A(X)\le t\}$, and let
$\mu_A=\Law(\ii_A(X))$. Define CDFs recursively by
\begin{equation}\label{eq:F}
F_0(t)=\mathbf1_{\{t\ge0\}},\qquad
F_n(t)=\min\{(\mu_P*F_{n-1})(t),(\mu_Q*F_{n-1})(t)\}.
\end{equation}
Convolution means $(\mu*F)(t)=\int F(t-x)\,\mu(dx)$.
The minimum of two finite-support CDFs is again a finite-support CDF.

\begin{lemma}[Profile lower bound and exact dyadic attainment]\label{lem:optimal}
For any finite $P,Q$, every admissible discrete horizon-$n$ seed $S$
satisfies $F_S(t)\le F_n(t)$ for all real $t$.
For pure dyadic $P,Q$, $F_n$ is the information CDF of a constructible
finite pure dyadic seed $R_n$ and
$C_n^{\mathrm{causal}}=H(R_n)$.
\end{lemma}
\begin{proof}
At horizon zero every seed has nonnegative information. Suppose the
profile bound holds for $n-1$ and fix the first action $a$ with output $Y$.
For every $y$ of positive probability, $S\mid Y=y$ is an admissible seed
for the remaining horizon. On that fiber,
\[
-\log_2\Prb\{S=s\}=
\ii_a(y)-\log_2\Prb\{S=s\mid Y=y\}.
\]
Induction gives $F_S(t)\le(\mu_a*F_{n-1})(t)$ for both first actions,
proving the lower bound. Integrating survival functions bounds seed
entropy below by the mean of $F_n$.

For finite pure dyadic laws $R,A$, a deterministic map from $R$ to $A$
exists if and only if $F_R(k)\le F_A(k)$ for all integers $k$.
Necessity follows because a source atom must fit inside its target atom.
For sufficiency, assign source atoms in decreasing order of size. At
mass $2^{-k}$, all eligible remaining target capacities are integer
multiples of $2^{-k}$ and total $F_A(k)-F_R(k-1)$, enough for the current
level. Thus all bins can be filled exactly.

For $F(k)=\min(F_A(k),F_B(k))$, the counts
$N_k=2^k(F(k)-F(k-1))$ are nonnegative integers. Taking $N_k$ atoms of
mass $2^{-k}$ realizes $F$ and refines both $A,B$.
Apply this recursively to $P\otimes R_{n-1}$ and $Q\otimes R_{n-1}$,
starting at a point mass $R_0$. The refinement maps send the current
state to an output and a new state. Conditional on that output, the new
state has exactly law $R_{n-1}$. Induction proves correctness after every
history and for every adaptive policy. The profile lower bound proves
optimality, even against non-dyadic seeds.
\end{proof}

\section{The envelope and its exact entropy increments}
Assume henceforth that both entropies equal $h$. Put
$\alpha=\Law(\ii_P-h)$ and $\beta=\Law(\ii_Q-h)$.
Both laws are centered. Let a controller choose either increment law
predictably, and let $S_j^\pi$ be its sum after $j$ steps. Define
\[
G_j(x)=\inf_\pi\Prb\{S_j^\pi\le x\},\qquad
U_j(x)=\sup_\pi\Prb\{S_j^\pi\le x\}.
\]
Dynamic programming, conditioning on the first increment, gives
\[
G_j=\min(\alpha*G_{j-1},\beta*G_{j-1}),\qquad
G_0=\mathbf1_{[0,\infty)}.
\]
Consequently $G_j(x)=F_j(jh+x)$.
Let $Z_j$ have CDF $G_j$ and put $e_j=\E Z_j$.
For all finite action laws $\Delta_j\ge e_j$; for pure dyadic laws,
Lemma~\ref{lem:optimal} gives equality.
The two CDFs in the minimum have the same mean $e_{j-1}$. Therefore
\begin{equation}\label{eq:increments}
e_{j+1}-e_j=\frac12\| (\alpha-\beta)*G_j\|_1\ge0.
\end{equation}
Indeed, use $\min(A,B)=(A+B-|A-B|)/2$ and the CDF formula for differences
of means. This identity charges a nonnegative cost at every horizon.

Write $B=\max\{|x|:x\in\operatorname{supp}\alpha\cup
\operatorname{supp}\beta\}$. If $B=0$, both information laws coincide
and the problem is trivial. Assume $B>0$. For every policy, conditional
centering and $|D_i|\le B$ imply
\[
\Prb\{S_j^\pi\ge t\},\ \Prb\{S_j^\pi\le-t\}
\le\exp\left(-\frac{t^2}{2jB^2}\right),\qquad t>0.
\]
For completeness, convexity gives
$\E[e^{sD_i}\mid\mathcal F_{i-1}]\le\cosh(sB)\le e^{s^2B^2/2}$;
iterate and optimize the exponential Markov bound.
The estimate is uniform over policies, hence also bounds both tails of
$Z_j$. Integrating the tails yields
\begin{equation}\label{eq:second}
\E Z_j^2\le4B^2j,\qquad 0\le e_j\le2B\sqrt j.
\end{equation}

\section{A Fourier lower bound: why a logarithm is unavoidable}
Let $D=\alpha-\beta$, and use the Fourier convention
$\widehat\nu(t)=\int e^{itx}\nu(dx)$.
Suppose $d\ge2$ is the first order at which the moments of $\alpha,\beta$
differ. Finite support and Taylor expansion imply that for some
$c_0,t_0>0$,
\begin{equation}\label{eq:fourierD}
|\widehat D(t)|\ge c_0|t|^d\qquad(0<|t|\le t_0).
\end{equation}
Choose $c=\min(t_0,1/(2B))$ and $t_j=c/\sqrt{j+1}$.
Equation~\eqref{eq:second} and $\cos u\ge1-u^2/2$ show
\[
|\widehat{\Law(Z_j)}(t_j)|\ge
\Re\widehat{\Law(Z_j)}(t_j)
\ge1-\tfrac12t_j^2\E Z_j^2\ge\tfrac12.
\]
The function $R_j=D*G_j$ is compactly supported, and its distributional
derivative is $D*\Law(Z_j)$. Integration by parts therefore gives
\[
\|R_j\|_1\ge|\widehat R_j(t_j)|
=\frac{|\widehat D(t_j)|\,
|\widehat{\Law(Z_j)}(t_j)|}{|t_j|}
\ge\frac{c_0c^{d-1}}2(j+1)^{-(d-1)/2}.
\]
Together with \eqref{eq:increments}, this proves
\begin{equation}\label{eq:lower}
e_n\ge\frac{c_0c^{d-1}}4
\sum_{j=0}^{n-1}(j+1)^{-(d-1)/2}.
\end{equation}
For $d=2$ this is $\Omega(\sqrt n)$; for $d=3$ it is
$\Omega(\log(n+1))$. Since $\Delta_n\ge e_n$ without a dyadic
assumption, Corollary~\ref{cor:general} follows already.
The upper estimate in \eqref{eq:second} proves the matching square-root
upper bound for pure dyadic laws.

\section{An integrated envelope contraction}
For the remaining upper bounds, put
\[
W_j=\int_{\mathbb R}(U_j-G_j)(x)\,dx.
\]
The all-$\beta$ walk has mean zero and its CDF lies between the envelopes.
Thus $0\le e_j\le W_j$. No reflection symmetry is assumed here.

Let $V$ be independent smoothing noise supported in $[-A,A]$, let $T_r$
be the all-$\beta$ walk, and define
\[
h_r(x)=\Prb\{T_r+V\le x\},\quad
R_r=D*h_r,\quad K_r=|R_r|,\quad
 a_r=\|K_r\|_1,\quad b_r=\|K_r'\|_1.
\]
The noise chosen below ensures that these functions are absolutely
continuous and compactly supported where needed.

\begin{lemma}[Contraction]\label{lem:contraction}
For every $n\ge1$,
\begin{equation}\label{eq:contract}
W_n\le2A+2\sum_{r=0}^{n-1}(a_r+b_rW_{n-1-r}).
\end{equation}
If $b_*:=\sum_{r\ge0}b_r<1/2$, then for every $N\ge1$,
\begin{equation}\label{eq:maxW}
\max_{0\le n\le N}W_n
\le\frac{2A+2\sum_{r=0}^{N-1}a_r}{1-2b_*}.
\end{equation}
\end{lemma}
\begin{proof}
Write $G_n^V,U_n^V$ for the smoothed policy envelopes and $W_n^V$ for
their integrated width. Bounded support of $V$ gives
$U_n(x)\le U_n^V(x+A)$ and $G_n(x)\ge G_n^V(x-A)$.
Every CDF $F$ satisfies $\int(F(x+A)-F(x-A))dx=2A$ by Tonelli, so
$W_n\le W_n^V+2A$.

For each policy, telescope $\E h_{n-i}(x-S_i^\pi)$ over $i$.
Since $h_{r+1}=\beta*h_r$, conditioning on the past gives
\[
\Prb\{S_n^\pi+V\le x\}-h_n(x)
=\sum_{i=1}^n\E\left[
\mathbf1_{\{\text{action }i\text{ uses }\alpha\}}
R_{n-i}(x-S_{i-1}^\pi)\right].
\]
Its absolute value is at most
$\sum_i\E K_{n-i}(x-S_{i-1}^\pi)$.
Let $B_j$ be the CDF of $T_j$. For every controlled prefix CDF $F_j^\pi$,
$|F_j^\pi-B_j|\le U_j-G_j$. Stieltjes integration by parts gives
\[
\E K_r(x-S_j^\pi)\le\E K_r(x-T_j)
+\int|K_r'(x-y)|(U_j(y)-G_j(y))\,dy.
\]
The right side is independent of $\pi$, so taking the supremum first
and then integrating in $x$ proves
\[
\int\sup_\pi\E K_r(x-S_j^\pi)dx\le a_r+b_rW_j.
\]
The smoothed envelope width is at most twice the largest absolute
deviation from $h_n$ at each threshold. Summing proves
\eqref{eq:contract}. Taking a maximum over $n\le N$, with $W_0=0$,
and moving $2b_*$ times that finite maximum to the left proves
\eqref{eq:maxW}.
\end{proof}

\section{Derivative estimates for an arbitrary finite reference law}
Assume $\beta$ is nondegenerate. Choose two of its support points
$x<y$, put $b=(y-x)/2>0$ and
$\rho=2\min(\beta\{x\},\beta\{y\})>0$.
Then $\beta=\rho\eta+(1-\rho)\zeta$, where
$\eta=(\delta_x+\delta_y)/2$ and $\zeta$ is a probability law when
$\rho<1$. The case $\rho=1$ needs no residual law.
Let $v_a$ denote the uniform density on $[-a,a]$, choose
$a=Lb$ with a positive integer $L$, and take $V$ to be the sum of
four independent variables with density $v_a$. Then $A=4a$ and, in
weak-derivative notation,
\begin{equation}\label{eq:derivative}
h_r^{(q+1)}=(v_a')^{*q}*v_a^{*(4-q)}*\beta^{*r},
\qquad q=2,3,4.
\end{equation}
Here $v_a'=(\delta_{-a}-\delta_a)/(2a)$, and exponent zero means
$\delta_0$. In particular $\|h_r^{(q+1)}\|_{\TV}\le a^{-q}$.
The fourth case can be a signed measure; for densities the total
variation norm agrees with the $L^1$ norm.

\begin{lemma}\label{lem:decay}
There are constants $C_q<\infty$ depending on $\beta,q$, but independent
of $a=Lb$ and $r$, such that
\begin{equation}\label{eq:decay}
\|h_r^{(q+1)}\|_{\TV}
\le\min\{a^{-q},C_q(r+1)^{-q/2}\},\qquad q=2,3,4.
\end{equation}
\end{lemma}
\begin{proof}
Telescoping point masses shows
$v_a'=\nu_L*v_b'$, where $\nu_L$ is uniform on
$\{-a+b,-a+3b,\ldots,a-b\}$.
Convolution with probability measures cannot increase total variation.
In the binomial mixture expansion of $\beta^{*r}$, the number $K$ of
$\eta$ increments has law $\operatorname{Bin}(r,\rho)$.
For a fixed $K=k$, translation and scaling of symmetric signs give
\[
\|(v_b')^{*q}*\eta^{*k}\|_{\TV}
\le b^{-q}\binom{k+q}{q}^{-1/2}.
\]
To verify this inequality, let $N=k+q$ independent symmetric signs be
$\varepsilon_i$. Up to sign and the factor $b^{-q}$, the signed measure
is the law of their sum weighted by $\prod_{i=1}^q\varepsilon_i$.
Given the sum, the conditional product equals the average of
$\prod_{i\in I}\varepsilon_i$ over all $q$-element subsets $I$.
These products are orthonormal, so that average has squared $L^2$ norm
$\binom Nq^{-1}$. Cauchy--Schwarz proves the asserted variation bound.

It remains to average over $K$. Using
$\binom{k+q}{q}\ge(k+1)^q/q!$ and splitting at $K=\rho r/2$ gives
\[
\E\binom{K+q}{q}^{-1/2}
\le\sqrt{q!}(1+\rho r/2)^{-q/2}+e^{-\rho r/8}
\le C'_q(r+1)^{-q/2}.
\]
For the exponential term, apply Markov to $e^{-(\ln2)K}$ and use
$1-\rho/2\le e^{-\rho/2}$ and $\ln2<3/4$.
This proves the decay estimate; combine it with the direct bound
$a^{-q}$ from \eqref{eq:derivative}.
\end{proof}

\section{Closing the classification}
For pure dyadic laws, the uncentered information laws have integer
support. If their moments through order $m-1$ coincide, their probability
generating polynomials differ by a multiple of $(1-z)^m$.
Centering by the common $h$ gives a finite signed measure $\Gamma$ such
that
\begin{equation}\label{eq:factor}
D=(\delta_0-\delta_1)^{*m}*\Gamma.
\end{equation}
This follows because ordinary moments and falling-factorial moments
span the same polynomials through that degree. The common translation
by $-h$ is absorbed into $\Gamma$; $h$ need not be an integer.

Let $M=\|\Gamma\|_{\TV}$. The $m$th unit finite difference is, up to
translation, the $m$th weak derivative convolved with
$\mathbf1_{[0,1]}^{*m}$, whose $L^1$ norm is one. Therefore
\begin{equation}\label{eq:ab}
a_r\le M\|h_r^{(m)}\|_{\TV},\qquad
b_r\le M\|h_r^{(m+1)}\|_{\TV}.
\end{equation}
If $v_P=v_Q>0$, the reference law is nondegenerate and
Lemma~\ref{lem:decay} applies.
For $m=3$ or $4$, the summable sequence
$MC_m(r+1)^{-m/2}$ dominates the bounds on $b_r$.
As $a=Lb\to\infty$, their pointwise upper bound $Ma^{-m}$ tends to
zero. Dominated convergence lets us fix a finite $a$ with
$\sum b_r<1/2$.

If variances match but third moments differ, use $m=3$.
Then $a_r\le MC_2/(r+1)$, and \eqref{eq:maxW} gives
$W_n=O(\log(n+1))$. Since $\Delta_n=e_n\le W_n$, the Fourier lower
bound proves the logarithmic case.

If variances and third moments both match, use $m=4$.
Then $a_r\le MC_3(r+1)^{-3/2}$ is summable, so
\eqref{eq:maxW} gives $W_n=O(1)$ and $\Delta_n=O(1)$.
If the common variance is zero, both information laws are the same point
mass and this conclusion holds trivially.
Finally \eqref{eq:increments} proves monotonicity. When
$\alpha\ne\beta$, it also gives
$e_1=\tfrac12\int|F_\alpha-F_\beta|>0$.
For pure dyadic laws, equality of the information laws is equivalent to
equality of all atom multiplicities, hence to being permutations of one
another. In that case the recursion is the ordinary product profile and
$\Delta_n=0$. This completes Theorem~\ref{thm:classification}.

\section{An explicit logarithmic pair, with numerical constants}
Take
\begin{align*}
P&=\bigl(1/16,(1/32)^{\times10},(1/64)^{\times12},
                       (1/128)^{\times56}\bigr),\\
Q&=\bigl((1/32)^{\times16},(1/128)^{\times64}\bigr).
\end{align*}
Their support sizes are $79$ and $80$. Their centered information laws are
\[
\alpha=\tfrac1{16}\delta_{-2}+\tfrac5{16}\delta_{-1}
       +\tfrac3{16}\delta_0+\tfrac7{16}\delta_1,
\qquad \beta=\tfrac12\delta_{-1}+\tfrac12\delta_1.
\]
Direct calculation gives
\[
h_P=h_Q=6,\qquad v_P=v_Q=1,\qquad
\kappa_P=-\tfrac38,\quad\kappa_Q=0,\qquad
D=\tfrac1{16}\delta_{-2}*(\delta_0-\delta_1)^{*3}.
\]
Let $\mathsf H_n=\sum_{k=1}^n1/k$ be the harmonic number.

\begin{theorem}[A genuine logarithmic overhead]\label{thm:explicit}
For the displayed fixed pair and every $n\ge1$,
\begin{equation}\label{eq:explicit}
\boxed{\frac{\mathsf H_n}{2048}
\le C_n^{\mathrm{causal}}(P,Q)-6n
<\frac{128}{5}+\frac3{10}\mathsf H_n.}
\end{equation}
In particular, the overhead is $\Theta(\log n)$ and is unbounded.
The exact one-step excess is $\Delta_1=1/8$.
\end{theorem}
\begin{proof}
For the lower bound, $B=2$ in \eqref{eq:second}. Choose
$t_j=1/(4\sqrt{j+1})$, so $|\widehat{\Law(Z_j)}(t_j)|\ge1/2$.
For $0\le t\le1/4$, the inequality
$2\sin(t/2)\ge t-t^3/24$ implies
$|1-e^{it}|^3\ge t^3/2$. Thus
$|\widehat D(t_j)|\ge t_j^3/32$, and
\eqref{eq:increments} gives
\[
e_{j+1}-e_j\ge\tfrac12\frac{t_j^3/32}{t_j}\cdot\tfrac12
=\frac1{2048(j+1)}.
\]
Summation proves the lower bound.

For the upper bound use four uniform smoothing variables on $[-2,2]$,
so $a=2$ and $A=8$. For the Rademacher reference, the derivative estimate
has the explicit form
\[
\|h_r^{(q+1)}\|_{\TV}
\le\min\{2^{-q},\binom{r+q}{q}^{-1/2}\}.
\]
Equations \eqref{eq:ab} with $m=3$ and $M=1/16$ give
\[
\sum_{r=0}^{n-1}a_r\le\frac{\sqrt2}{16}\mathsf H_n,
\qquad
b_*\le\frac1{16}\sum_{r\ge0}
\min\{2^{-3},\sqrt6(r+1)^{-3/2}\}
\le\frac{3\sqrt[3]6}{32}.
\]
For the last sum, bound it by
$\int_0^\infty\min(2^{-3},\sqrt6 x^{-3/2})dx
=3\sqrt[3]6/2$.
Therefore \eqref{eq:maxW} and $e_n\le W_n$ give
\[
\Delta_n\le
\frac{16+\frac{\sqrt2}{8}\mathsf H_n}
     {1-\frac{3\sqrt[3]6}{16}}
<\frac{128}{5}+\frac3{10}\mathsf H_n,
\]
using $\sqrt2<3/2$ and $\sqrt[3]6<2$.
Finally the CDF of $\alpha-\beta$ is $1/16,-2/16,1/16$ on the unit
intervals starting at $-2,-1,0$. Equation~\eqref{eq:increments} gives
$\Delta_1=(1+2+1)/32=1/8$.
\end{proof}

\section{Scope and reproducibility}
The classification is for all finite pure dyadic pairs of equal entropy,
with no symmetry assumption and no restriction on their alphabet sizes.
For general finite laws the lower bounds remain valid, but this note
does not supply matching upper constructions. It does not claim a full
classification outside the pure dyadic class or literature priority.

The companion program \texttt{logarithmic\_seed.py} constructs the exact
profile of the explicit pair, verifies normalization and entropy using
integer atom counts, and checks complete decoder partitions through
horizon two. It also checks the moment identities, cubic factorization,
and displayed harmonic bounds at finite horizons. Optional floating-point
runs illustrate growth but are not part of the all-horizon proof.
For example, a floating-point run gives
$\Delta_{1000}\simeq0.62701621$ and
$\Delta_{10000}\simeq0.69669408$.
Slow apparent stabilization does not contradict logarithmic divergence.

Each horizon uses its own finite initial seed, with no later randomness.
Expanded state spaces may be exponential in the horizon. No single
projectively consistent infinite source or efficient expanded decoder is
asserted. The proof is independent of numerical extrapolation.
\end{document}
